% 第 0 章：集合、逻辑演算与证明武器库
% 说明：elegantbook 的 cn 方案会自动给 \chapter 加“第 X 章”前缀，
\chapter{集合、逻辑演算与证明武器库}

\section{等价变形：防增根与防漏根}

\begin{problem}\level{易}
解方程 $\sqrt{x+3}=x+1$，并说明为什么两边平方后必须检验。
\end{problem}
\begin{solution}
\paragraph{集合与逻辑结构剖析}
平方不是同解变形：由 $a=b$ 可推出 $a^2=b^2$，反之不然，因为 $a^2=b^2\Leftrightarrow a=b$ 或 $a=-b$。于是 $\sqrt{x+3}=x+1$ 与平方后所得方程不同解，后者的解集可能更大，必须代回原方程检验。
\paragraph{标准解题步骤}
原方程要求 $x+3\geq0$ 且右边 $x+1\geq0$（算术平方根非负），即 $x\geq-1$。平方得 $x+3=(x+1)^2=x^2+2x+1$，即 $x^2+x-2=0$，$(x+2)(x-1)=0$，解得 $x=-2$ 或 $x=1$。检验：$x=-2$ 时 $x+1=-1<0$，左边 $\sqrt{x+3}=1$，两边不等，舍去；$x=1$ 时左边 $\sqrt4=2$，右边 $2$，成立。故解集为 $\{1\}$。
\paragraph{学生典型易错点诊断}
最易犯的错误是平方后忘记检验，把增根 $x=-2$ 当作解；二是只写 $x+3\geq0$ 而漏掉 $x+1\geq0$，定义域判断不全；三是误以为“两边平方”永远等价；四是不检验就下结论。规避方法：先求原方程有意义的范围，平方后解方程，再逐一回代检验。
\end{solution}

\begin{problem}\level{中}
解关于 $x$ 的方程 $(a-1)x=a^2-1$，并讨论参数 $a$ 取不同值时方程解的情况。
\end{problem}
\begin{solution}
\paragraph{集合与逻辑结构剖析}
两边同除以含参因式 $a-1$，相当于默认 $a-1\neq0$。除法只在除式非零时保持同解；当 $a-1=0$ 时原方程退化为 $0\cdot x=0$，解集完全不同。故含参因式必须先按“是否为零”分类，不能一除了之。
\paragraph{标准解题步骤}
原方程即 $(a-1)x=(a-1)(a+1)$。当 $a-1\neq0$ 即 $a\neq1$ 时，两边同除以 $a-1$ 得 $x=a+1$，方程有唯一解；当 $a-1=0$ 即 $a=1$ 时，方程化为 $0\cdot x=0$，任意实数 $x$ 都是解，解集为 $\mathbb{R}$。综上：$a\neq1$ 时解集为 $\{a+1\}$；$a=1$ 时解集为 $\mathbb{R}$。
\paragraph{学生典型易错点诊断}
最易犯的错误是不讨论 $a-1$ 是否为零就两边同除，漏掉 $a=1$ 时无穷多解的情形；二是把 $a=1$ 误判为无解，混淆“恒等”与“矛盾”；三是把 $x=a+1$ 当成对所有 $a$ 成立，忽略前提 $a\neq1$；四是求出 $x=a+1$ 后不回代检验。规避方法：见含参因式先设其不为零再除，再单独讨论其为零的情形。
\end{solution}

\section{逻辑演算：德摩根律与逆否等价}

\begin{problem}\level{中}
设集合 $A=\{x\mid x^2-5x+6\leq0\}$，$B=\{x\mid x^2-2ax+a^2-1\leq0\}$。若 $A\cap B$ 中“至多含有一个整数”，求实数 $a$ 的取值范围。
\end{problem}
\begin{solution}
\paragraph{集合与逻辑结构剖析}
“至多含有一个整数”正面分类较繁（$0$ 个、$1$ 个），用德摩根律取反面更简：其否定是“至少含两个整数”。而 $A=[2,3]$ 中整数只有 $2,3$，故反面即“$2$ 与 $3$ 同时属于 $B$”。求出反面范围再取补集即可。
\paragraph{标准解题步骤}
由 $x^2-5x+6\leq0$ 得 $(x-2)(x-3)\leq0$，即 $A=[2,3]$，其中整数仅为 $2,3$。由 $x^2-2ax+a^2-1\leq0$ 得 $(x-a)^2\leq1$，即 $B=[a-1,a+1]$。反面“至少两个整数”等价于 $2,3\in B$。$2\in B\iff a-1\leq2\leq a+1\iff1\leq a\leq3$；$3\in B\iff a-1\leq3\leq a+1\iff2\leq a\leq4$。两者同时成立时 $a\in[1,3]\cap[2,4]=[2,3]$。由德摩根律，所求为其补集：$a<2$ 或 $a>3$，即 $(-\infty,2)\cup(3,+\infty)$。
\paragraph{学生典型易错点诊断}
最易犯的错误是正面按“$0$ 个或 $1$ 个整数”分类，情形多且易漏；二是忘记 $A$ 中只有 $2,3$ 两个整数，误以为整数有无穷多个；三是把“至多一个”的否定写成“至少一个”（应为“至少两个”）；四是求出 $[1,3]\cap[2,4]=[2,3]$ 后忘记取补集，直接写成答案。规避方法：见“至多/至少”先写否定，用德摩根律化“且”为“或”，求出反面再整体取补。
\end{solution}

\begin{problem}\level{中}
设 $p:|x-1|\leq2$，$q:x^2-2x+1-m^2\leq0\ (m>0)$。若 $\neg p$ 是 $\neg q$ 的充分不必要条件，求实数 $m$ 的取值范围。
\end{problem}
\begin{solution}
\paragraph{集合与逻辑结构剖析}
核心工具是逆否命题等价：$\neg p\Rightarrow\neg q$ 与 $q\Rightarrow p$ 同真同假。设 $P=\{x\mid p\text{ 成立}\}$，$Q=\{x\mid q\text{ 成立}\}$，则“$\neg p$ 是 $\neg q$ 的充分条件”等价于 $Q\subseteq P$；“不必要”等价于 $P\not\subseteq Q$，合起来即 $Q\subsetneq P$。
\paragraph{标准解题步骤}
由 $|x-1|\leq2$ 得 $-1\leq x\leq3$，即 $P=[-1,3]$。由 $x^2-2x+1-m^2\leq0$ 得 $(x-1)^2\leq m^2$，又 $m>0$，故 $1-m\leq x\leq1+m$，即 $Q=[1-m,1+m]$。由逆否命题等价，$\neg p$ 是 $\neg q$ 的充分不必要条件 $\iff Q\subsetneq P$。$Q\subseteq P\iff\begin{cases}1-m\geq-1,\\1+m\leq3,\end{cases}\iff m\leq2$。当 $m=2$ 时 $Q=[-1,3]=P$，不满足真包含，故 $0<m<2$。
\paragraph{学生典型易错点诊断}
最易犯的错误是把“$\neg p$ 是 $\neg q$ 的充分条件”直接翻译成 $P\subseteq Q$，不用逆否命题转换，方向出错；二是忽略“不必要”对应真包含，漏掉 $m=2$ 的检验，错答 $0<m\leq2$；三是丢掉 $m>0$ 的条件，写出负数区间；四是解 $q$ 时只写上界、漏掉下界。规避方法：先用逆否命题把 $\neg p,\neg q$“翻”成 $q,p$，再按集合包含关系列不等式，最后检验等号是否使两集合相等。
\end{solution}

\section{量词否定与间接证明}

\begin{problem}\level{难}
设 $f(x)=x^2-2x+3$，$x\in[0,3]$。
(1) 若对任意 $x\in[0,3]$，$f(x)\geq a$ 恒成立，求 $a$ 的取值范围；
(2) 若存在 $x\in[0,3]$，使 $f(x)\geq a$ 成立，求 $a$ 的取值范围；
(3) 写出“$\exists x\in[0,3],\ f(x)\geq a$”的否定，并求使它为真时 $a$ 的取值范围。
\end{problem}
\begin{solution}
\paragraph{集合与逻辑结构剖析}
量词决定取最值的方向：$\forall x,\ f(x)\geq a$ 恒成立 $\iff a$ 不超过 $f$ 的最小值（求最小）；$\exists x,\ f(x)\geq a$ 能成立 $\iff a$ 不超过 $f$ 的最大值（求最大）。量词否定须量词与结论同时否定：$\neg(\exists x, f(x)\geq a)=\forall x, f(x)<a$。
\paragraph{标准解题步骤}
$f(x)=x^2-2x+3=(x-1)^2+2$，对称轴 $x=1\in[0,3]$，故最小值 $f(1)=2$；端点 $f(0)=3$，$f(3)=6$，故最大值 $6$。(1) 恒成立 $\iff a\leq f_{\min}=2$。(2) 能成立 $\iff a\leq f_{\max}=6$。(3) 否定为“$\forall x\in[0,3],\ f(x)<a$”，它为真 $\iff a>f_{\max}=6$。
\paragraph{学生典型易错点诊断}
最易犯的错误是把“$\forall$ 恒成立”误求最大值，错答 $a\leq6$；二是把“$\exists$ 能成立”误求最小值，错答 $a\leq2$；三是写否定时只改量词不改结论（或只改结论不改量词）；四是忽略严格不等号，把 (3) 错答 $a\geq6$。规避方法：先分离参数，再按“任意取最小、存在取最大”定方向，否定时量词与结论同时否定。
\end{solution}

\begin{problem}\level{难}
用反证法证明：$\sqrt{2}$ 不是有理数。要求写出“假设—推导—矛盾”的完整链条，并指出矛盾具体出现在哪一步。
\end{problem}
\begin{solution}
\paragraph{集合与逻辑结构剖析}
反证法的逻辑骨架是：要证 $p$，先假设 $\neg p$，由 $\neg p$ 推出与已知公理、定义或已证结论相矛盾的结论，从而断定 $\neg p$ 不成立，故 $p$ 成立。其依据是逆否命题等价：$\neg p\Rightarrow q$ 且 $\neg q$ 为真，则 $p$ 为真。
\paragraph{标准解题步骤}
假设 $\sqrt2$ 是有理数，则存在互质的正整数 $p,q$ 使 $\sqrt2=\dfrac pq$。两边平方得 $p^2=2q^2$，故 $p^2$ 为偶数，从而 $p$ 为偶数，设 $p=2k$。代入得 $4k^2=2q^2$，即 $q^2=2k^2$，故 $q^2$ 为偶数，$q$ 也为偶数。于是 $p,q$ 同为偶数、有公因数 $2$，与“$p,q$ 互质”矛盾。故假设不成立，$\sqrt2$ 不是有理数。
\paragraph{学生典型易错点诊断}
最易犯的错误是只假设“$\sqrt2$ 是有理数”而不设 $p,q$ 互质，导致推不出矛盾；二是不用“$p^2$ 偶则 $p$ 偶”这一关键引理，推导断裂；三是把矛盾说成“$p,q$ 都是偶数”本身，其实矛盾是“与互质假设冲突”，要落到假设上；四是靠 $\sqrt2\approx1.414\ldots$ 下结论，这不是证明。规避方法：反证法三步走——假设结论反面、推出矛盾、肯定结论，并明确“与什么矛盾”。
\end{solution}

\section{数学归纳法与补集思想}

\begin{problem}\level{中}
用数学归纳法证明：当 $n\in\mathbb{N}$ 且 $n\geq5$ 时，$2^n>n^2$。
\end{problem}
\begin{solution}
\paragraph{集合与逻辑结构剖析}
数学归纳法是“从有限到无限”的桥梁，两步缺一不可：归纳基础（验证首项）与归纳递推（由 $n=k$ 成立推 $n=k+1$ 成立）。本题递推的关键是证明 $2k^2\geq(k+1)^2$，它只在 $k$ 足够大时成立，故必须从 $n=5$ 起步。
\paragraph{标准解题步骤}
(1) 基础：$n=5$ 时 $2^5=32>25=5^2$，成立。(2) 递推：设 $n=k\ (k\geq5)$ 时 $2^k>k^2$ 成立，则 $n=k+1$ 时 $2^{k+1}=2\cdot2^k>2k^2$。要证 $2k^2\geq(k+1)^2$，即 $k^2-2k-1\geq0$，即 $(k-1)^2\geq2$；由 $k\geq5$ 得 $(k-1)^2\geq16>2$，故 $2k^2>(k+1)^2$，从而 $2^{k+1}>(k+1)^2$。由数学归纳法，对一切 $n\geq5$，$2^n>n^2$。
\paragraph{学生典型易错点诊断}
最易犯的错误是只验证 $n=1$ 就宣布成立，但 $n=2,3,4$ 时并不成立（如 $2^3=8<9=3^2$），归纳基础必须选 $n=5$；二是递推时直接写 $2^{k+1}>2k^2>(k+1)^2$，默认中间不等式而不加证明；三是忘记说明 $k\geq5$ 这一前提的作用；四是用 $n=k$ 的结论去证 $n=k$（循环论证）。规避方法：先找最小起始值并验证，再明确“假设 $n=k$ 成立”去推 $n=k+1$，并单独证明递推中需要的中间不等式。
\end{solution}

\begin{problem}\level{易}
已知关于 $x$ 的方程 $|x-1|+|x-3|=a$ 无实数解，求实数 $a$ 的取值范围。
\end{problem}
\begin{solution}
\paragraph{集合与逻辑结构剖析}
“无实数解”的正面讨论（按绝对值零点分段）较繁，其反面“有实数解”等价于“$a$ 落在函数 $f(x)=|x-1|+|x-3|$ 的值域中”，求值域反而简单。这正是补集思想：先求反面，再在参数全集 $\mathbb{R}$ 中取补。
\paragraph{标准解题步骤}
设 $f(x)=|x-1|+|x-3|$。由绝对值三角不等式 $f(x)\geq|(x-1)-(x-3)|=2$，当 $1\leq x\leq3$ 时取等号，故 $f$ 的值域为 $[2,+\infty)$。方程 $f(x)=a$ 有实数解 $\iff a\in[2,+\infty)\iff a\geq2$。取补集，方程无实数解 $\iff a<2$。
\paragraph{学生典型易错点诊断}
最易犯的错误是正面按 $x<1$、$1\leq x\leq3$、$x>3$ 三段讨论，计算量大且易错；二是忘记取补集，直接把“有解”的范围 $a\geq2$ 当成答案；三是误用绝对值不等式方向，写成 $f(x)\leq2$；四是把“无解”等价于“判别式小于零”，但本题不是二次方程，没有判别式可用。规避方法：见“无解/不存在”先想反面“有解/存在”，求出反面范围后务必在参数全集内取补。
\end{solution}
